% Formalization support lemmas.
%
% NOT book material. These declarations were invented in this workspace to
% carry a Lean formalization and have no counterpart in the source text --
% they carry no \dcref. They live here, outside the book chapters, so the
% book blueprint stays timeless mathematics while their \leanok marks and
% Lean anchors are preserved.
%
% Do not add book mathematics to this file, and do not add formalization
% notes to the book chapters.

\begin{remark}
\label{topping-ch10-perelman-reference}
\dcref{10.2.3}
\group{positive-ricci-three-manifolds}
\source{topping:page-0127}
This severely constrains the singularities which are possible. To understand
more about the Ricci flow beyond these notes, the reader is directed first to
Perelman's paper.
\end{remark}

\begin{remark}
\label{topping-ch8-classical-entropy-link}
\dcref{8.2.4}
\group{perelman-w-entropy}
\source{topping:page-0092}
The $\mathcal W$-functional applied to this backwards Brownian diffusion on a
Ricci flow also arises via the renormalised classical entropy $\tilde N$.
\end{remark}

\begin{remark}
\label{topping-ch8-no-collapsing-perelman}
\dcref{8.3.3}
\group{perelman-w-entropy}
\source{topping:page-0096}
This is modelled on Perelman's no collapsing theorem from [31, \S4].
\end{remark}

\begin{remark}
\label{topping-ch8-no-collapsing-variant}
\dcref{8.3.2}
\group{perelman-w-entropy}
\source{topping:page-0096}
We could replace the $|R| \le r^{-2}$ hypothesis by $|R| \le C r^{-2}$,
although $\xi$ would then also depend on $C$.
\end{remark}

\begin{remark}
\label{topping-gradient-flow-volume-preserving-remark}
\dcref{6.4.2}
\group{topping-gradient-flow}
\source{topping:page-0074}
In the construction above, the measure $\omega$ can be described either from the initial data $(\hat g(0),\hat f^{(0)})$ or from the terminal data $(\hat g(T),f(T))$ after pulling back by $\psi_T$.
\end{remark}

\begin{example}
\label{topping-intro-cigar-soliton}
\dcref{Hamilton's cigar soliton}
\group{ricci-solitons}
\source{topping:page-0011}
On $\mathbb R^2$, let $g_0=\rho^2(dx^2+dy^2)$ with
\[
\rho^2=\frac{1}{1+x^2+y^2}.
\]
Then
\[
\Ric(g_0)=\frac{2}{1+x^2+y^2}g_0,
\]
and for
\[
Y=-2\!\left(x\frac{\partial}{\partial x}+y\frac{\partial}{\partial y}\right)
\]
one has
\[
\mathcal L_Y g_0=-\frac{4}{1+x^2+y^2}g_0.
\]
Hence $g_0$ is a steady Ricci soliton. In geodesic polar coordinates
\[
g_0=ds^2+\tanh^2(s)\,d\theta^2,
\]
with Gauss curvature
\[
K=\frac{2}{\cosh^2 s}.
\]
It is also a gradient soliton, with potential function
\[
f=-2\ln(\cosh s).
\]
\end{example}

\begin{remark}
\label{topping-maxpr-curvature-doubling}
\dcref{3.2.12}
\group{maximum-principle}
\source{topping:page-0050}
If $|\Rm|\le 1$ at $t=0$, then the waiting time before the maximum of $|\Rm|$
doubles is bounded below by a positive constant depending only on the dimension
$n$.
\end{remark}

\begin{remark}
\label{topping-maxpr-perelman-volume}
\dcref{3.2.8}
\group{maximum-principle}
\source{topping:page-0048}
If the Ricci flow exists for all $t\in[0,\infty)$, then
\[
\overline V:=\lim_{t\to\infty}\frac{V(t)}{\left(1+\frac{2(-\alpha)}{n}t\right)^{\frac n2}}
\]
exists. A rough principle in Perelman's work is that, locally, $\overline V>0$
suggests the manifold is becoming hyperbolic, while $\overline V=0$ suggests a
graph-manifold-like limit.
\end{remark}

\begin{remark}
\label{topping-maxpr-preserved-curvature-conditions}
\dcref{3.2.9}
\group{maximum-principle}
\source{topping:page-0048}
Scalar curvature is only one curvature quantity whose positivity is preserved
under Ricci flow on closed manifolds. Other preserved conditions include
positive curvature operator in any dimension, positive isotropic curvature in
four dimensions, and positive Ricci curvature in dimension three or less.
Negative curvatures are not preserved in general.
\end{remark}

\begin{remark}
\label{topping-pde-principal-symbol-fourier-convention}
\dcref{4.2.1}
\group{topping-pde-existence}
\source{topping:page-0056}
Some authors write $i\xi$ where this text writes $\xi$, in accordance with the
Fourier transform convention.
\end{remark}
